\begin{figure}[H]
\centering
\begin{tikzpicture}[
    node distance=1.1cm, 
    artifact/.style={rectangle, draw, fill=blue!5, font=\footnotesize, inner sep=4pt, align=center},
    process/.style={rectangle, draw, fill=orange!5, font=\footnotesize\itshape, inner sep=4pt, align=center},
    verdict/.style={rectangle, draw, fill=green!10, rounded corners, font=\footnotesize\bfseries, inner sep=4pt, align=center},
    arrow/.style={thick, ->, >=stealth}
]

% --- FIRST LINE ---
\node (bytecode) [artifact] {.class\\bytecode};
\node (soot)     [artifact, right=of bytecode, xshift=0.4cm] {Jimple};
\node (callgen)  [artifact, right=of soot, xshift=0.5cm, yshift=0.8cm] {Call Graph};
\node (overridegen) [artifact, right=of soot,xshift=0.3cm, yshift=-0.8cm] {Override Graph};
\node (cg)       [artifact, right=of soot, xshift=3.6cm, ] {Merged Graph};
\node (jdk_cache) [artifact, above=of callgen, fill=gray!10] {JDK Summary\\Cache};

% --- SECOND LINE (Positioned below the CallGraph) ---
\node (loop)     [artifact, right=of cg, xshift=1.2cm] {Batch Order};
\node (ptg)      [artifact, below=of loop, yshift=-0.5cm] {Points-to\\Graph and\\Write Set};

% Final split outputs
\node (verdict)  [artifact,  left=of ptg, xshift=-4.5cm] {Verdict};
\node (summary)  [artifact, left=of verdict, xshift=-1.0cm, fill=gray!10] {MethodSummary\\(Local Cached)};

% --- PATHS ---
% Top row connections
\draw [arrow] (bytecode) -- node[above, font=\scriptsize\ttfamily] {SootUp} (soot);
\draw [arrow] (soot) -- (callgen);
\draw [arrow] (soot) -- (overridegen);
\draw [arrow] (jdk_cache.south) -- (callgen.north);
\draw [arrow] (callgen) -- (cg);
\draw [arrow] (overridegen) -- (cg);

% Connection between rows (The "Snake" turn)
\draw [arrow] (cg) -- node[above, font=\scriptsize\ttfamily, align=center] {Tarjan's SCC}(loop);

% Bottom row connections (Moving right to left)
\draw [arrow] (loop) -- node[above, font=\scriptsize\ttfamily, align=center] {dataflow analysis}  (ptg);
\draw [arrow] (ptg) -- node[above, font=\scriptsize\ttfamily, align=center] {side effect inference} (verdict);

% Caching connection
\draw [arrow] (verdict) -- (summary);
\draw [arrow, dashed] (summary.west) 
    -- ++(-2.3,0) 
    -- ++(0,5.9) 
    -- ([yshift=0.3cm]jdk_cache.north) 
    node[midway, above, font=\small] {if is JDK method} 
    -- (jdk_cache.north);

% Feedback loop back to the loop node for next method
\draw [arrow, dashed, gray] (verdict.north) -- (loop.south)
    node[pos=0.5, above, font=\tiny, black] {iterate};

\end{tikzpicture}
\caption{Flow Chart of Implementation}
\label{fig:implement-flowchart}
\end{figure}

\paragraph{JDK summary cache.}
Analyzing user code inevitably requires resolving calls into the JDK standard library.
Rather than re-analyzing library methods on every invocation, we persist method summaries to a disk-backed JSON cache (\texttt{jdk-cache/}).
When a cross-file callee is encountered during call graph construction or dataflow analysis, the tool first consults this cache; only if no summary exists does it perform on-demand analysis of the library method.
We ship a pre-computed cache covering common \texttt{java.util}, \texttt{java.lang}, and \texttt{java.io} methods, so that users benefit from library-level precision out of the box without incurring the cost of whole-JDK analysis.